\documentclass[11pt,reqno]{amsart}

\usepackage[T1]{fontenc}
\usepackage{lmodern}
\usepackage{amsmath,amssymb,mathtools}
\usepackage{booktabs}
\usepackage{enumitem}
\usepackage{microtype}
\usepackage[colorlinks=true,linkcolor=blue,citecolor=blue,urlcolor=blue]{hyperref}

\newtheorem{theorem}{Theorem}[section]
\newtheorem{proposition}[theorem]{Proposition}
\newtheorem{lemma}[theorem]{Lemma}
\newtheorem{corollary}[theorem]{Corollary}
\theoremstyle{definition}
\newtheorem{definition}[theorem]{Definition}
\newtheorem{example}[theorem]{Example}
\theoremstyle{remark}
\newtheorem{remark}[theorem]{Remark}

\DeclareMathOperator{\adj}{adj}
\DeclareMathOperator{\rem}{rem}
\DeclareMathOperator{\rank}{rank}
\DeclareMathOperator{\tr}{tr}
\DeclareMathOperator{\Spec}{Spec}
\DeclareMathOperator{\cont}{cont}
\DeclareMathOperator{\N}{N}
\DeclareMathOperator{\disc}{disc}
\newcommand{\Q}{\mathbf Q}
\newcommand{\Z}{\mathbf Z}
\newcommand{\R}{\mathbf R}
\newcommand{\Gm}{\mathbf G_m}

\title[Spectral arithmetic surgery]{Spectral Arithmetic Surgery:\\
Pinned Full--Minor Assignment, All--Rank Arithmetic Obstructions,\\
and Exact Salem Height Surgery}
\author{KEIO}
\date{September 2026}

\begin{document}

\begin{abstract}
Let $T$ be a rational matrix self-adjoint for a positive rational form $H$,
let $e=e_1$ be cyclic, and prescribe simultaneously the characteristic
polynomial of $T$ and that of its first principal compression.  We develop an
exact arithmetic surgery calculus for low-rank $H$-self-adjoint
perturbations and several rank-free obstructions tied to a fixed rational
symmetrizing form.  In rank one, a prescribed rational full--minor pair is
attainable if and only if its cross polynomial is a scalar square of a mixed
cofactor.  This yields a finite cyclic-coset integrality test and a sharp
quartic Salem obstruction for every $K\geq5$, with an exception at $K=4$.
In rank two we derive the exact Gram--Pl\"ucker law, including its exterior
correction, prove the converse assignment theorem, and reduce every fixed
two-plane problem to rational linear elimination and one quadric.

For the canonical order-three realization we prove two all-rank barriers.
Semidefinite output-pin surgery collapses to rank one, while arbitrary
rational surgery preserving the canonical form must satisfy an explicit
binary norm equation.  A local obstruction to this equation produces an
infinite family of genuine height-one Salem targets unattainable at every
rank.  On the canonical rank-two plane we strengthen the exact square-class
criterion to a rational parametrization and finite integer-square sieve and
show that every strict branch is necessarily indefinite.  The unique
integral target capable of a vanishing exterior shift is birational to the
rank-zero elliptic curve of Cremona label $240a3$; its known rational points,
also recovered by a short $2$-isogeny descent, exclude every $K\geq5$.
In contrast, a sporadic $K=5$ certificate and an infinite family
$K=\nu^2+1$ give exact indefinite rank-two height collapse.  These results
close several natural surgery routes but do not prove a Salem-number closure
conjecture.
\end{abstract}

\maketitle

\section{Introduction}

The characteristic polynomial of a rank-one perturbation is governed by the
matrix determinant lemma.  The inverse problem becomes subtler when one asks
to prescribe, at the same time, the characteristic polynomial of a fixed
principal compression.  It becomes arithmetic when both outputs must be
integral and geometrically constrained by interlacing.  The purpose of this
paper is to isolate the exact algebraic image of this simultaneous operation
in ranks one and two and to test it on a rigid family of Salem pencils.

The classical ingredients---adjugates, resolvents, cyclic vectors, low-rank
determinant formulas, restricted-rank spectral modification, principal-minor
relations, and interlacing---are standard; see, for example,
\cite{dptz,merzel,merzelm,kula,arbenzgangergolub,
alahmadiehvinzant,delrio}.  The contribution here is their simultaneous
full--first-principal-compression arithmetic packaging for one fixed cyclic
model and the resulting exact classification and obstruction theorems.  The
Salem quartics used below are classical specializations of Boyd's criterion
\cite{boydquartic,zaimi}; no novelty is claimed for the polynomials
themselves.

Here is the main chain of results.

\begin{enumerate}[label=\textup{(\roman*)}]
\item A rational output pair $(f_*,P_*)$ is produced by a fixed cyclic
rank-one model if and only if
\[
 fP_*-f_*P=tA^2
\]
for an explicitly reconstructible mixed cofactor $A$.
\item Two output pins force $X^2(X-2)^2$ to divide this cross polynomial.
For a fixed pin direction, simultaneous integrality is membership in one
finite cyclic subgroup of $\Q[X]/\Z[X]$.
\item The quartics
\[
 z^4-(K-1)z^3+z^2-(K-1)z+1
\]
have exact normal height $K$, but for $K\geq5$ their canonical pencils admit
no nonzero signed two-pin rank-one surgery with integral Salem output.
\item In rank two the square is replaced by a binary Gram term minus an
exterior correction.  This identity again has a converse, hence it is an
exact assignment criterion rather than merely a necessary condition.
For a fixed update plane, feasibility reduces further to rational linear
elimination and one quadric.
\item Fixing the full and compressed eigenvectors is much stronger than
fixing the output roots.  An effective rank-$r$ update fixing a full
eigenline and a $d$-dimensional compressed eigenspace satisfies
$r\leq N-1-d$.  In the strict height-one slice, semidefinite output pins force
these eigenspace pins, so every-rank semidefinite surgery is excluded for
$K\geq5$.
\item For the canonical rational symmetrizing form, every rational target of
any rank satisfies an explicit binary norm equation.  Primes
$p\equiv5,7\pmod8$ give an infinite family of height-one Salem targets for
which this equation is locally insoluble.
\item For the plane generated by $T_Ke,T_K^2e$, the complete rational image
of the height-one target slice is a conic.  We parametrize it, obtain a finite
integer-square sieve at fixed $K$, and prove that every strict branch is
indefinite.
\item The sole integral target capable of a vanishing exterior shift is
birational to the elliptic curve $240a3$.  Its eight rational points exclude
all $K\geq5$ and imply a uniform asymptotic gap of at least $1/2$ from the
base exterior trace throughout the attainable integral slice.
\item In the opposite direction, the slice contains a sporadic $K=5$
certificate and an infinite integral Salem family with base height
$K=\nu^2+1$ and output height one.  None of these outputs occurs in rank one.
\end{enumerate}

The positive and negative results delimit one concrete surgery program.
Height collapse is possible, but in the complete slice it is necessarily
indefinite; the semidefinite, fixed-eigenspace, same-form local-obstruction,
and canonical vanishing-shift routes are closed.  Changing the symmetrizing
form, changing the update plane, or moving to higher order remains outside
these obstructions.  In particular, no statement below proves a Boyd closure
conjecture.

\section{Cyclic self-adjoint pencils and cofactors}

Throughout, unless stated otherwise,
\begin{equation}\label{eq:base}
 T\in M_N(\Q),\qquad H=\operatorname{diag}(1,H')\in M_N(\Q),\quad H>0,
 \qquad HT=T^{\mathsf T}H .
\end{equation}
Put $e=e_1$, let $B=T^{(1)}$ be the first principal compression, and set
\begin{equation}\label{eq:fP}
 f=\det(XI-T),\qquad P=\det(XI-B).
\end{equation}
We assume that $e$ is cyclic for $T$ and, when interlacing is invoked, that
$f$ and $P$ are separable and coprime.  Positive $H$-self-adjointness then
places the problem in the usual real self-adjoint category after conjugating
by $H^{1/2}$.

\begin{lemma}[Rational symmetrizable realization]\label{lem:realization}
Let $f,P\in\Q[X]$ be monic of degrees $N,N-1$, respectively, with simple
real zeros in strict interlacing position.  Then there is a rational
tridiagonal matrix
\[
 T=\begin{pmatrix}
 a_0&1&&\\ b_1&a_1&\ddots&\\ &\ddots&\ddots&1\\&&b_{N-1}&a_{N-1}
 \end{pmatrix},\qquad b_i\in\Q_{>0},
\]
whose characteristic pair is $(f,P)$.  It is self-adjoint for
\[
 H=\operatorname{diag}\bigl(1,b_1^{-1},(b_1b_2)^{-1},\ldots,
 (b_1\cdots b_{N-1})^{-1}\bigr),
\]
and $e_1$ is cyclic.
\end{lemma}

\begin{proof}
The rational function $P/f$ has simple positive residues at the zeros of
$f$.  Its finite Stieltjes--Jacobi continued fraction is obtained by the
Euclidean algorithm over $\Q[X]$.  Strict interlacing makes every
off-diagonal product $b_i$ positive, while monicity normalizes each
superdiagonal entry to one.  The continuant recurrence gives the stated
characteristic pair.  The displayed diagonal form then satisfies
$HT=T^{\mathsf T}H$, and nonzero off-diagonal products make $e_1$ cyclic;
compare \cite{hochstadt,deboor}.
\end{proof}

For $u=(u_1,u')\in\Q^N$ define
\begin{align}
 A_u(X)&=e^{\mathsf T}\adj(XI-T)u,\label{eq:Adef}\\
 C_u(X)&=u^{\mathsf T}H\adj(XI-T)u,\label{eq:Cdef}\\
 D_u(X)&=(u')^{\mathsf T}H'\adj(XI-B)u'.\label{eq:Ddef}
\end{align}

\begin{lemma}[Cyclic cofactor normal form]\label{lem:rankone-normal}
For every $u\in\Q^N$ there is a unique $p\in\Q[X]$, $\deg p<N$, with
$u=p(T)e$.  For this polynomial,
\begin{equation}\label{eq:normalform}
 A_u=\rem_f(pP),\qquad C_u=\rem_f(p^2P),
 \qquad PC_u-A_u^2=fD_u.
\end{equation}
\end{lemma}

\begin{proof}
Cyclicity gives the unique representation.  If $\lambda$ is a zero of $f$,
the simple-root adjugate formula, or equivalently spectral resolution of the
resolvent, gives
\[
 A_u(\lambda)=p(\lambda)P(\lambda),\qquad
 C_u(\lambda)=p(\lambda)^2P(\lambda).
\]
The degree bounds yield the two remainder formulas.  Write
$G=(XI-T)^{-1}$.  Self-adjointness gives
\[
 e^{\mathsf T}Gu=u^{\mathsf T}HGe=A_u/f,\qquad
 u^{\mathsf T}HGu=C_u/f.
\]
Apply Sherman--Morrison to $T+tuu^{\mathsf T}H$.  Comparing the $(1,1)$
resolvent entry, which equals the compression determinant divided by the full
determinant, gives
\[
 \frac{P_t}{f_t}=\frac Pf+
 \frac{t(A_u/f)^2}{1-tC_u/f}.
\]
Together with $f_t=f-tC_u$ this is precisely
$P_t=P-t(PC_u-A_u^2)/f$.  The second determinant lemma applied to the
compression identifies the quotient with $D_u$.
\end{proof}

\section{The rank-one square law and its converse}

For $t\in\Q$ consider the signed $H$-self-adjoint update
\begin{equation}\label{eq:rankoneupdate}
 T_t=T+tuu^{\mathsf T}H.
\end{equation}
The determinant lemma and Lemma~\ref{lem:rankone-normal} give
\begin{equation}\label{eq:rankoneoutputs}
 f_t=f-tC_u,\qquad P_t=P-tD_u,\qquad
 \boxed{fP_t-f_tP=tA_u^2}.
\end{equation}

\begin{theorem}[Exact rank-one full--minor assignment]\label{thm:rankone-iff}
Fix the cyclic realization \eqref{eq:base}--\eqref{eq:fP}, and let
$f_*,P_*$ be monic rational polynomials of degrees $N,N-1$.  For a fixed
$t\in\Q^\times$, the following are equivalent.
\begin{enumerate}[label=\textup{(\alph*)}]
\item There is $u\in\Q^N$ such that $T+tuu^{\mathsf T}H$ has full and first
compressed characteristic polynomials $f_*$ and $P_*$.
\item There is $A\in\Q[X]$, $\deg A<N$, such that
\begin{equation}\label{eq:squareiff}
 fP_*-f_*P=tA^2.
\end{equation}
\end{enumerate}
In case \textup{(b)}, the update vector is reconstructed uniquely from
\begin{equation}\label{eq:reconstructp}
 pP\equiv A\pmod f,\qquad \deg p<N,\qquad u=p(T)e.
\end{equation}
\end{theorem}

\begin{proof}
Necessity is \eqref{eq:rankoneoutputs}.  Conversely, $P$ is a unit in
$\Q[X]/(f)$, so \eqref{eq:reconstructp} has a unique solution.  With the
resulting $u$, reduce \eqref{eq:squareiff} modulo $f$ and use
Lemma~\ref{lem:rankone-normal}; this gives $tC_u=f-f_*$ as an equality because
both sides have degree $<N$.  Hence the full output is $f_*$.  Substituting
in $PC_u-A_u^2=fD_u$ and in \eqref{eq:squareiff} gives
$tD_u=P-P_*$, so the compressed output is $P_*$.
\end{proof}

\subsection{Two pins}

Assume that
\begin{equation}\label{eq:pins}
 f(0)=0,\quad P(0)\ne0,\qquad P(2)=0,\quad f(2)\ne0,
\end{equation}
and that the displayed roots are simple.  If a nonzero rank-one move retains
the output roots $f_t(0)=0$ and $P_t(2)=0$, then
$C_u(0)=0$, hence $p(0)=A_u(0)=0$.  Evaluating the cofactor identity at $2$
and using $D_u(2)=0$ gives $A_u(2)=0$.  Therefore
\begin{equation}\label{eq:pin-square-div}
 X(X-2)\mid A_u,\qquad
 X^2(X-2)^2\mid fP_t-f_tP.
\end{equation}
When $t$ may vary, the quotient in \eqref{eq:pin-square-div} must be a
nonzero rational scalar times a square.  By Theorem~\ref{thm:rankone-iff},
this square-class filter is also sufficient for the fixed realization.

\section{Arithmetic in a fixed pin direction}

Let $\phi_0$ span the zero eigenspace of $T$, and let $\psi_2$ span the
$2$-eigenspace of $B$.  Suppose the two rational linear constraints below are
independent and put
\begin{equation}\label{eq:pinspace}
 W=\{u:\langle u,\phi_0\rangle_H=0,
           \ \langle u',\psi_2\rangle_{H'}=0\}.
\end{equation}
Then $\dim W=N-2$.  For $u\in W$ write
\begin{equation}\label{eq:EG}
 C_u=XE_u,\qquad D_u=(X-2)G_u,\qquad
 f=Xr,\quad P=(X-2)q.
\end{equation}

\begin{proposition}[Cyclic-coset integrality criterion]\label{prop:cycliccoset}
Assume $f=Xr$ with $r\in\Z[X]$, write
$P=(X-2)q$ with $q\in\Q[X]$, and fix nonzero $u\in W$.  Then
\begin{equation}\label{eq:integralitycriterion}
 f_t,P_t\in\Z[X]\quad\Longleftrightarrow\quad
 tE_u\in\Z[X]\ \text{ and }\ q-tG_u\in\Z[X].
\end{equation}
For a nonzero rational polynomial $E=\sum_i e_iX^i$, define
$\cont_{\Q}(E)>0$ by
\[
 \sum_i\Z e_i=\cont_{\Q}(E)\Z,
\]
and set
$\alpha_u=\cont_{\Q}(E_u)^{-1}$.  Then
\begin{equation}\label{eq:cycliccoset}
 \text{\eqref{eq:integralitycriterion} is solvable with $t\ne0$}
 \quad\Longleftrightarrow\quad
 [q]\in\langle[\alpha_uG_u]\rangle
 \ \text{in }\Q[X]/\Z[X].
\end{equation}
\end{proposition}

\begin{proof}
The output formulas and \eqref{eq:EG} give
$f_t=X(r-tE_u)$ and $P_t=(X-2)(q-tG_u)$, proving
\eqref{eq:integralitycriterion}.  By the definition of rational content,
$tE_u$ is integral exactly when $t=n\alpha_u$ for some $n\in\Z$.  Substitution
gives \eqref{eq:cycliccoset}.  Since the coefficients are rational, the
generated subgroup is finite.  Consequently, if a class is represented only
by $n=0$ in a first choice of representative, adding the positive order of
$[\alpha_uG_u]$ produces a nonzero representing integer and hence $t\ne0$.
\end{proof}

Thus one direction is decided by Smith normal form or by a one-variable
Chinese-remainder calculation.  There is also a uniform denominator
statement once directions are normalized.

\begin{proposition}[Uniform content on a fixed pencil]\label{prop:uniformcontent}
Let $N\geq3$, let $\Lambda=W\cap\Z^N$, and choose an integral basis $S$ of
$\Lambda$.  There are positive integers $L,\Delta$, depending only on the
fixed pencil and this basis, such that for every primitive
$z\in\Z^{N-2}$,
\begin{equation}\label{eq:contentbound}
 \cont_{\Z}(LE_{Sz})\mid\Delta.
\end{equation}
Consequently the denominators of scalars producing integral full output are
uniformly bounded on primitive integral directions.
\end{proposition}

\begin{proof}
Choose $L$ so that the coefficients of $LE_{Sz}$ are integral homogeneous
quadrics $Q_1(z),\ldots,Q_m(z)$.  They have no common projective zero over
$\overline\Q$: otherwise $C_{Sz}=0$, and the simple-spectrum adjugate formula
\[
 C_{Sz}(\theta)=f'(\theta)
 \frac{\langle Sz,\varphi_\theta\rangle_H^2}
      {\langle\varphi_\theta,\varphi_\theta\rangle_H}
\]
would make $Sz$ orthogonal to every eigenvector, hence zero.  The projective
Nullstellensatz supplies an integral B\'ezout certificate.  Evaluating it at
a primitive $z$ shows that the gcd of all $Q_i(z)$ divides a fixed integer
$\Delta$, which is \eqref{eq:contentbound}.  The denominator assertion follows
from Proposition~\ref{prop:cycliccoset}.
\end{proof}

The primitive normalization is essential: replacing $u$ by $mu$ and $t$ by
$t/m^2$ represents the same matrix update.  The proof also isolates a local
obstruction.  If at a prime $p$ the coefficient quadrics of $E$ have no
projective base point modulo $p$, then some coefficient of $E_{Sz}$ is a
$p$-unit for every primitive $z$.  Hence $tE_{Sz}$ integral forces
$v_p(t)\geq0$; if $G$ is integral at $p$, such a scalar cannot cancel a
nonintegral coefficient of $q$.  We call this a denominator-prime base-locus
obstruction, not a ramification theorem.

\section{An infinite quartic rank-one obstruction}

For $K\geq3$ set
\begin{equation}\label{eq:rK}
 r_K=X^2-(K-1)X-1,\quad f_K=Xr_K,\quad
 P_K=(X-2)\left(X+\frac1K\right).
\end{equation}
The reciprocal lift is
\begin{equation}\label{eq:BK}
 B_K(z)=z^2r_K(z+z^{-1})
 =z^4-(K-1)z^3+z^2-(K-1)z+1.
\end{equation}
This is a genuine quartic Salem polynomial.  Indeed $r_K$ has roots
$y_K>2$ and $\alpha_K\in(-2,0)$ with $y_K\alpha_K=-1$ and
$K-1<y_K<K$; its discriminant lies strictly between consecutive squares.
The signs of $y_K^2-4$ at the two real embeddings then show that the
reciprocal lift is irreducible.  The same conclusion follows directly from
the classical quartic criterion \cite{zaimi}.

\begin{definition}[Normal height]
If a Salem polynomial of degree $2s$ is written
$R(z)=z^sr(z+z^{-1})$ and $f_R=Xr$, a strict integral normal interlacer is a
polynomial $b\in\Z[X]$ of degree $s-1$, with negative leading coefficient and
one simple zero in every gap between the compact roots of $f_R$.  Its normal
height is
\[
 \kappa(R)=\min_b\bigl(-\operatorname{lc}b\bigr).
\]
For the quartic $B_K$, this says precisely that $b=-\kappa X+c$ is primitive
and its zero lies strictly between $\alpha_K$ and $0$.
The corresponding monic compressed polynomial is
\begin{equation}\label{eq:normalbridge}
 P_b=-\frac{(X-2)b}{\kappa}.
\end{equation}
\end{definition}

\begin{proposition}[Exact height]\label{prop:exactheight}
For every $K\geq3$, $\kappa(B_K)=K$.
\end{proposition}

\begin{proof}
The zero $c/\kappa$ lies in $(\alpha_K,0)=(-1/y_K,0)$, so $c<0$ and
$|c|/\kappa<1/y_K$.  Primitivity gives $|c|\geq1$, while $y_K>K-1$;
hence $\kappa\geq K$.  The polynomial $-KX-1$ has zero $-1/K$ in the
required interval and attains the bound.
\end{proof}

\begin{theorem}[Quartic rank-one obstruction]\label{thm:quarticobstruction}
For every $K\geq5$, no nonzero signed rational rank-one update of the fixed
pair \eqref{eq:rK} can simultaneously
\begin{enumerate}[label=\textup{(\roman*)}]
\item preserve the output roots $0$ and $2$,
\item make both output polynomials integral, and
\item have the trace-root signature of a genuine quartic Salem polynomial.
\end{enumerate}
The threshold is sharp: $K=4$ admits such an output.
\end{theorem}

\begin{proof}
Write an integral pinned target as
\[
 f_*=Xr_*,\quad r_*=X^2+aX+b,\qquad
 P_*=(X-2)(X+n),\qquad a,b,n\in\Z.
\]
Theorem~\ref{thm:rankone-iff} and \eqref{eq:pin-square-div} reduce the cross
identity to
\begin{equation}\label{eq:quarticcross}
 r_K(X+n)-r_*\left(X+\frac1K\right)=\mu X(X-2)
\end{equation}
for some $\mu\in\Q^\times$.  Coefficient comparison gives
\begin{equation}\label{eq:quarticcoeffs}
 b=-Kn,\qquad
 a=\frac{3Kn-2K^2+K-2}{2K+1},\qquad
 \mu=-\frac{(K-1)(Kn-1)}{K(2K+1)}.
\end{equation}
Since $2K+1$ is odd, $a\in\Z$ is equivalent to
\begin{equation}\label{eq:congruence}
 2K+1\mid3(n+2).
\end{equation}
Put $g=\gcd(3,2K+1)$ and $L=(2K+1)/g$.  Thus $n=-2+L\ell$.

If $\ell=0$, then $r_*=(X-2)(X-K)$, so the root $2$ collides.  If
$\ell\leq-1$, evaluation of \eqref{eq:quarticcross} at $2$ gives
\[
 r_*(2)=\frac{(5-2K)(n+2)}{2+1/K}>0,
\]
whereas a monic quadratic having one root in $(-2,2)$ and one root above $2$
must be negative at $2$.  Finally, if $\ell\geq1$, then $n\geq3$ for
$K\geq5$, and direct substitution gives
\[
 r_*(-2)=\frac{4K^2+6K+8-Kn(2K+7)}{2K+1}<0.
\]
The desired two roots both lie to the right of $-2$, which would require
$r_*(-2)>0$.  All cases are excluded.

For $K=4$, take
\[
 n=1,\qquad r_*=X^2-2X-4,\qquad \mu=-\frac14.
\]
Its roots are $1\pm\sqrt5$, one in $(-2,2)$ and the other above $2$.
Moreover $r_*$ is irreducible because its discriminant is $20$, and
$r_*(2)r_*(-2)=-16<0$ shows that $y^2-4$ is not a square in its trace
field.  Hence the reciprocal lift is an irreducible quartic Salem polynomial.
Equation \eqref{eq:quarticcross} and Theorem~\ref{thm:rankone-iff} supply the
actual signed rational update.
\end{proof}

The obstruction is not a failure of irreducibility or of a thin-set
argument.  The integral lattice itself misses the open Salem chamber.  It
therefore cannot be removed by strong approximation or Hilbert
irreducibility.

\section{Rank-two Gram--Pl\"ucker surgery}

Let $U=(u_1,u_2)\in M_{N,2}(\Q)$, let $U'$ be obtained by deleting its first
row, and let $S=S^{\mathsf T}\in M_2(\Q)$.  Define
\begin{align}
 A&=e^{\mathsf T}\adj(XI-T)U\in M_{1,2}(\Q[X]),\label{eq:matrixA}\\
 C&=U^{\mathsf T}H\adj(XI-T)U\in\operatorname{Sym}_2(\Q[X]),\label{eq:matrixC}\\
 D&=(U')^{\mathsf T}H'\adj(XI-B)U'
 \in\operatorname{Sym}_2(\Q[X]).\label{eq:matrixD}
\end{align}
Polarization of Lemma~\ref{lem:rankone-normal} gives
\begin{equation}\label{eq:matrixcofactor}
 \boxed{PC-A^{\mathsf T}A=fD}.
\end{equation}
At each simple zero of $f$, the matrix $C$ has rank at most one, hence
$f\mid\det C$; similarly $P\mid\det D$.  Put
\begin{equation}\label{eq:phis}
 \Phi_C=\frac{\det C}{f},\qquad
 \Phi_D=\frac{\det D}{P}.
\end{equation}
Taking determinants in \eqref{eq:matrixcofactor} yields the polynomial
\begin{equation}\label{eq:Psi}
 \boxed{\Psi_U=\frac{A\adj(C)A^{\mathsf T}}{f}
 =P\Phi_C-f\Phi_D\in\Q[X]}.
\end{equation}

For the update
\begin{equation}\label{eq:ranktwoupdate}
 T_S=T+USU^{\mathsf T}H,
\end{equation}
the $2\times2$ determinant lemma, valid also for singular $S$, gives
\begin{equation}\label{eq:ranktwooutputs}
 f_S=f-\tr(SC)+\det(S)\Phi_C,\qquad
 P_S=P-\tr(SD)+\det(S)\Phi_D.
\end{equation}

\begin{theorem}[Gram--Pl\"ucker law]\label{thm:GP}
The cross polynomial of \eqref{eq:ranktwoupdate} is
\begin{equation}\label{eq:GP}
 \boxed{fP_S-f_SP=ASA^{\mathsf T}-\det(S)\Psi_U.}
\end{equation}
The product $\det(S)\Psi_U$ is invariant under changing the basis of the
update plane.
\end{theorem}

\begin{proof}
Substitute \eqref{eq:ranktwooutputs}, use
$PC-fD=A^{\mathsf T}A$, and use
$P\Phi_C-f\Phi_D=\Psi_U$.  Under
$U\mapsto UL$ and $S\mapsto L^{-1}SL^{-\mathsf T}$,
$\Psi_U$ is multiplied by $\det(L)^2$ and $\det S$ by $\det(L)^{-2}$.
\end{proof}

When $S=tcc^{\mathsf T}$ has rank one, the determinant term vanishes and
\eqref{eq:GP} is the square law.  For nonsingular $S$, rational congruence
diagonalization expresses the first term as a binary Gram form, but the
nonzero exterior correction remains.  This is why the rank-two image is not
described by a bare binary norm form.

\begin{theorem}[Exact rank-at-most-two assignment]\label{thm:ranktwoiff}
Let $f_*,P_*$ be monic rational polynomials of degrees $N,N-1$.  The following
are equivalent.
\begin{enumerate}[label=\textup{(\alph*)}]
\item The pair $(f_*,P_*)$ is the full--minor output of
$T+USU^{\mathsf T}H$ for rational $U\in M_{N,2}(\Q)$ and symmetric rational
$S$.
\item There are $A_1,A_2\in\Q[X]_{<N}$ and $S=S^{\mathsf T}\in M_2(\Q)$
such that, after uniquely reconstructing
\[
 p_iP\equiv A_i\pmod f,\quad
 C_{ij}=\rem_f(p_ip_jP),\quad
 D=\frac{PC-A^{\mathsf T}A}{f},
 \]
and then setting
\[
 \Phi_C=\frac{\det C}{f},\qquad
 \Phi_D=\frac{\det D}{P},\qquad
 \Psi=P\Phi_C-f\Phi_D,
\]
one has
\begin{equation}\label{eq:ranktwoiff}
 fP_*-f_*P=ASA^{\mathsf T}-\det(S)\Psi.
\end{equation}
\end{enumerate}
The update has rank exactly two when $p_1,p_2$ are independent modulo $f$
and $\det S\ne0$.
\end{theorem}

\begin{proof}
Necessity is Theorem~\ref{thm:GP}.  Conversely set $u_i=p_i(T)e$ and let
$(\widetilde f,\widetilde P)$ be the actual output.  Equations
\eqref{eq:GP} and \eqref{eq:ranktwoiff} imply
\[
 f(\widetilde P-P_*)=P(\widetilde f-f_*).
\]
Since $(f,P)=1$, the polynomial $f$ divides $\widetilde f-f_*$.  The latter
has degree at most $N-1$, so it vanishes.  The compressed difference then
vanishes as well.  The final assertion follows because $H$ is nonsingular,
the columns of $U$ are independent, and $S$ is nonsingular.
\end{proof}

\subsection{Fixed-plane feasibility as one quadric}

The converse above also turns fixed-plane feasibility into an exact
local--global calculation rather than a search through update matrices.

\begin{proposition}[Fixed-plane linear--quadric reduction]
\label{prop:fixed-plane-quadric}
Fix a rational two-plane $U$ and write
\[
 S=\begin{pmatrix}r&s\\s&u\end{pmatrix},\qquad d=\det S.
\]
For prescribed rational monic polynomials $f_*,P_*$, coefficient matching in
\eqref{eq:ranktwooutputs} is a rational affine linear system
\begin{equation}\label{eq:linearquadric}
 \mathcal L_{U,f_*,P_*}(r,s,u,d)=b.
\end{equation}
The pair $(f_*,P_*)$ is attainable in this plane if and only if
\eqref{eq:linearquadric} has a rational solution satisfying
\begin{equation}\label{eq:detquadric}
 d=ru-s^2.
\end{equation}
After rationally parametrizing the affine solution space of
\eqref{eq:linearquadric}, feasibility is therefore the existence of a
rational point on one explicit affine quadric.

If the projective closure of the residual quadric is nonsingular and has
positive dimension, feasibility is equivalent to feasibility over $\R$ and
over every $\Q_p$.  Singular and zero-dimensional residual problems reduce,
by rational linear algebra, to a lower-rank quadratic, a discriminant test,
or a linear equation.
\end{proposition}

\begin{proof}
The first assertion is immediate from
\[
 f_S=f-\tr(SC)+d\Phi_C,\qquad
 P_S=P-\tr(SD)+d\Phi_D:
\]
all coefficients are affine linear in $(r,s,u,d)$, and the only condition
lost by treating $d$ as an independent variable is
\eqref{eq:detquadric}.  Conversely, a common solution defines a symmetric
$S$ with determinant $d$, and \eqref{eq:ranktwooutputs} gives the prescribed
pair.

Gaussian elimination parametrizes every nonempty rational affine solution
space.  Substitution in \eqref{eq:detquadric} leaves one quadratic equation.
For a nonsingular projective closure of positive dimension, the
Hasse--Minkowski theorem gives a rational projective point exactly when there
is a point over every completion; weak approximation for quadrics then moves
a rational point away from the hyperplane at infinity.  The degenerate cases
are obtained by splitting off the radical and solving the remaining
quadratic or linear equation; see \cite{cassels}.
\end{proof}

\section{Output pins versus fixed eigenspaces}

If only $f_S(0)=0$ and $P_S(2)=0$ are required, then tautologically
\begin{equation}\label{eq:ranktwo-output-pin}
 X(X-2)\mid fP_S-f_SP.
\end{equation}
For indefinite $S$, cancellation can preserve an eigenvalue while changing
its eigendirection.  Thus \eqref{eq:ranktwo-output-pin} need not contain
squared pin factors.

Pointwise eigenspace preservation is different.  The following theorem is
stated over a general field because it is linear algebra rather than
positivity.

\begin{lemma}[Active congruence factorization]\label{lem:activefactor}
Let the ground field have characteristic different from $2$, let
$H=H^{\mathsf T}$ be nonsingular, and suppose
$H\Delta=\Delta^{\mathsf T}H$.  If $r=\rank\Delta$, then
\[
 \Delta=USU^{\mathsf T}H
\]
for a full-column-rank $N\times r$ matrix $U$ and a nonsingular symmetric
$r\times r$ matrix $S$.
\end{lemma}

\begin{proof}
The matrix $M=\Delta H^{-1}$ is symmetric of rank $r$.  A symmetric matrix
over a field of characteristic different from $2$ has a rank factorization
$M=USU^{\mathsf T}$ obtained by restricting its bilinear form to a
nondegenerate complement of its radical.  Here $U$ has full column rank and
$S$ is symmetric and nonsingular.  Multiplication by $H$ gives the claim.
\end{proof}

\begin{theorem}[All-rank fixed-eigenspace barrier]\label{thm:allrankbarrier}
Let the ground field have characteristic different from $2$.  Suppose
$H=H^{\mathsf T}=\operatorname{diag}(h,H')$ is nonsingular and
$HT=T^{\mathsf T}H$.
Let
\[
 T\phi_0=\lambda_0\phi_0,\qquad B\psi=\lambda_2\psi,\qquad
 e^{\mathsf T}H\phi_0\ne0.
\]
Let $\Delta=USU^{\mathsf T}H$ have effective rank $r$, with the inactive
null directions removed so that $U$ has full column rank and $S$ is
nonsingular.  If $T+\Delta$ fixes $\phi_0$ with eigenvalue $\lambda_0$ and
its first compression fixes $\psi$ with eigenvalue $\lambda_2$, then
\begin{equation}\label{eq:rankbarrier}
 r\leq N-2.
\end{equation}
More generally, if a $d$-dimensional eigenspace of $B$ is fixed pointwise,
then $r\leq N-1-d$.
\end{theorem}

\begin{proof}
The full pin gives
$0=USU^{\mathsf T}H\phi_0$, hence $U^{\mathsf T}H\phi_0=0$ by injectivity of
$U$ and $S$.  Deletion of the first coordinate is injective on
$\operatorname{im}U$.  Indeed, if $Uc=ae$, then
\[
 0=(Uc)^{\mathsf T}H\phi_0=a e^{\mathsf T}H\phi_0,
\]
so $a=0$ and then $c=0$.  Thus the deleted matrix $U'$ still has rank $r$.
The compressed pin gives $(U')^{\mathsf T}H'\psi=0$, placing
$\operatorname{im}U'$ in one hyperplane of an $(N-1)$-dimensional space.
This proves \eqref{eq:rankbarrier}.  For a pointwise fixed $d$-space, the
image lies in its $H'$-orthogonal complement, of dimension $N-1-d$.
\end{proof}

\begin{corollary}[Order-three collapse]\label{cor:orderthreecollapse}
For $N=3$, every finite-rank $H$-self-adjoint perturbation fixing the full
eigenline and the compressed eigenline as above has effective rank at most
one.  Hence Theorem~\ref{thm:quarticobstruction} survives every nominal-rank
attempt under pointwise fixed-eigenspace pins.
\end{corollary}

\begin{proof}
Apply Lemma~\ref{lem:activefactor} and
Theorem~\ref{thm:allrankbarrier} with $d=1$.
\end{proof}

There is one important situation in which output pins automatically become
eigenspace pins.  For a positive form $H$, call an $H$-self-adjoint
perturbation $\Delta$ $H$-positive (respectively $H$-negative) semidefinite
when the symmetric matrix $H\Delta$ is positive (respectively negative)
semidefinite.

\begin{lemma}[Semidefinite equality forces an eigenline pin]
\label{lem:semidefinite-pin}
Let $A=A^{\mathsf T}$ and $E=E^{\mathsf T}$ be real matrices with
$E\succeq0$ or $E\preceq0$.  Suppose that the ordered eigenvalue
$\lambda_j(A)$ is simple and
\[
 \lambda_j(A+E)=\lambda_j(A).
\]
Then $E$ annihilates an eigenvector of $A$ belonging to $\lambda_j(A)$.
\end{lemma}

\begin{proof}
Assume first that $E\succeq0$.  The ordered eigenvalue
$\lambda_j(A+tE)$ is nondecreasing for $0\leq t\leq1$; see, for example,
\cite{bhatia}.  Equality at the endpoints makes it constant.  If $v$ is a
unit eigenvector for the simple initial eigenvalue, the
Hellmann--Feynman formula gives
\[
 \left.\frac{d}{dt}\lambda_j(A+tE)\right|_{t=0}=v^{\mathsf T}Ev=0.
\]
Positive semidefiniteness implies $Ev=0$.  The negative-semidefinite case
follows by applying the same argument to $-E$.
\end{proof}

\begin{theorem}[All-rank semidefinite obstruction in the strict slice]
\label{thm:allrank-semidefinite}
Let $K\geq5$ and let $\Delta$ be a rational $H$-self-adjoint perturbation of
a positive rational realization of the pair \eqref{eq:rK}.  Suppose that
$H\Delta$ is semidefinite and that its output pair is
\begin{equation}\label{eq:semidefinite-target}
\begin{gathered}
 f_*=X(X^2-aX-K),\qquad P_*=(X-2)(X+1),\\
 a\in\Z,\qquad \frac{K-4}{2}<a<K-1.
\end{gathered}
\end{equation}
Then no such perturbation exists.
\end{theorem}

\begin{proof}
Conjugate by $H^{1/2}$.  At both endpoints, $0$ is the simple middle full
eigenvalue: the other two full roots consist of one point in $(-2,-1)$ and
one point above $2$.  Likewise, $2$ is the simple maximal eigenvalue of both
compressed matrices.  Lemma~\ref{lem:semidefinite-pin}, applied to the full
matrix and to its principal compression, shows that $\Delta$ fixes the
initial full zero eigenline and the initial compressed $2$-eigenline
pointwise.  Theorem~\ref{thm:allrankbarrier} now gives
$\rank\Delta\leq1$.  The zero update cannot change
$P_K=(X-2)(X+1/K)$ to $P_*$, while every nonzero rank-one update is excluded
by Theorem~\ref{thm:quarticobstruction}.
\end{proof}

\begin{remark}\label{rem:semidefinite-sharp}
The threshold is sharp.  The $K=4$ exception in the proof of
Theorem~\ref{thm:quarticobstruction} has square-law scalar $-1/4$ and hence
is a negative-semidefinite rational rank-one update.  The theorem above is
not a claim about arbitrary semidefinite perturbations: the integral output
pair, the two ordered pins, and the strict Salem signature in
\eqref{eq:semidefinite-target} are essential.
\end{remark}

In the positive rational setting of \eqref{eq:pins}, take an active rank-two
representation.  Pointwise preservation is columnwise orthogonality, so
the spectral adjugate formula gives $A_i(0)=0$.  At the compressed pin one
has $D(2)=0$; evaluating \eqref{eq:matrixcofactor} at $2$ gives
$A(2)^{\mathsf T}A(2)=0$, hence $A_i(2)=0$ over $\Q$.  Thus every $A_i$ is
divisible by $\eta=X(X-2)$.  Since $A_i(0)=p_i(0)P(0)$ and $P(0)\ne0$, the full
pin also gives $p_i(0)=0$; hence every entry of $C$ is divisible by $X$ and, for a
$2\times2$ matrix, so is every entry of $\adj(C)$.  Writing $f=Xr$ gives
\[
 \Psi=\frac{A\adj(C)A^{\mathsf T}}{f}
      =\eta^2\frac{N(X)}{r(X)}
\]
for some $N\in\Q[X]$.  Since $\Psi$ is a polynomial and
$\gcd(r,\eta)=1$, one has $r\mid N$, and therefore $\eta^2\mid\Psi$.
Thus both terms in \eqref{eq:GP} carry the squared pin divisor.  The
constructions below deliberately require only output pins; they fix the zero
eigenline but change the compressed $2$-eigenline.

\section{A rank-free obstruction with the canonical form}

We now fix more structure.  Put
\begin{equation}\label{eq:canonical-b}
 \mathcal B_K=2K^3-6K^2+3K-1,\qquad
 b_1=\frac{\mathcal B_K}{K^2},\qquad
 b_2=\frac{2K(K-1)(2K-5)}{\mathcal B_K^2},
\end{equation}
and let $T_K$ be the rational Jacobi realization of $(f_K,P_K)$ displayed
in Appendix~\ref{app:explicit}.  Its positive symmetrizing form is
\begin{equation}\label{eq:canonical-H}
 H_K=\operatorname{diag}\bigl(1,b_1^{-1},(b_1b_2)^{-1}\bigr).
\end{equation}
The next obstruction is independent of an update plane and of the effective
rank, but it is tied to this fixed form and to the first-coordinate
compression.

\begin{theorem}[Rank-free same-form norm obstruction]
\label{thm:sameH-norm}
Let $a,m\in\Q$.  Suppose that a rational $H_K$-self-adjoint matrix $T_*$
has characteristic pair
\begin{equation}\label{eq:sameH-target}
 f_*=X(X^2-aX-m),\qquad P_*=(X-2)(X+1).
\end{equation}
Then the affine conic
\begin{equation}\label{eq:sameH-norm}
 Y^2+2K(K-1)(2K-5)Z^2
   =K^2\mathcal B_K(a+m-3)
\end{equation}
has a rational point.  Consequently, \eqref{eq:sameH-norm} is necessary for
a rational $H_K$-self-adjoint perturbation of $T_K$ of any effective rank.
\end{theorem}

\begin{proof}
Every rational $H_K$-self-adjoint matrix has the form
\[
 T_*=\begin{pmatrix}
 x&y&z\\
 b_1y&u&v\\
 b_1b_2z&b_2v&w
 \end{pmatrix}.
\]
The compressed characteristic polynomial in \eqref{eq:sameH-target} gives
\[
 u+w=1,\qquad uw-b_2v^2=-2,
\]
while the full trace gives $x=a-1$.  Comparing the coefficient of $X$ in
$\det(XI-T_*)=X^3-aX^2-mX$ yields
\[
 b_1y^2+b_1b_2z^2=a+m-3.
\]
Multiplication by $K^2\mathcal B_K$ and the substitution
$Y=\mathcal B_Ky$, $Z=z$ give \eqref{eq:sameH-norm}.

The conclusion does not depend on a presentation of the perturbation.  The
difference of two rational $H_K$-self-adjoint matrices is again
$H_K$-self-adjoint, and Lemma~\ref{lem:activefactor} represents it as
$USU^{\mathsf T}H_K$ with effective rank equal to its matrix rank.
\end{proof}

\begin{corollary}[An infinite fixed-form, all-rank Salem obstruction]
\label{cor:allrank-Salem-family}
Let $p>5$ be prime with
\[
 p\equiv5\ \text{or}\ 7\pmod8,\qquad K=\frac{p+5}{2},
\]
and set $a=K-4$, $m=K$.  The pair
\begin{equation}\label{eq:allrank-target}
 f_*=X\bigl(X^2-(K-4)X-K\bigr),\qquad
 P_*=(X-2)(X+1)
\end{equation}
is a genuine height-one Salem target, but it cannot be reached from the
fixed canonical realization $(T_K,H_K)$ by a rational
$H_K$-self-adjoint perturbation of any rank.
\end{corollary}

\begin{proof}
Here $2K-5=p$, $a+m-3=p-2$, and
$\mathcal B_K\equiv1/4\pmod p$.  The right side of
\eqref{eq:sameH-norm} is therefore a $p$-adic unit whose residue square class
is
\[
 K^2\mathcal B_K(p-2)\equiv-\frac{25}{8}
 =-2\left(\frac54\right)^2\pmod p.
\]
Since $(-2/p)=-1$ for $p\equiv5,7\pmod8$, the conic has no $\Q_p$-point.
Indeed, in an equation $Y^2+pcZ^2=d$ with $c,d\in\Z_p^\times$, the two
summands have valuations of opposite parity.  A negative valuation of either
variable would give the left side negative valuation, so $Y,Z\in\Z_p$;
reduction modulo $p$ would then make the nonsquare $d$ a square.

It remains to check the target.  For
$q=X^2-(K-4)X-K$ one has
\[
 q(-2)=K-4>0,\qquad q(-1)=-3<0,\qquad
 q(2)=-3(K-4)<0.
\]
Thus the roots lie in $(-2,-1)$ and $(2,\infty)$.  Its discriminant is
$(K-2)^2+12$, which is not a square for $K\geq6$: a square would give
$(s-K+2)(s+K-2)=12$, whose only positive same-parity factorization leads to
$K=4$.  Moreover
\[
 q(2)q(-2)=-3(K-4)^2<0,
\]
so $y^2-4$ is not a square in the quadratic trace field.  Hence the
reciprocal lift is irreducible and Salem.  The primitive normal interlacer
$-(X+1)$ gives height one.  The obstruction now follows from
Theorem~\ref{thm:sameH-norm}.
\end{proof}

The restriction to the canonical form is essential.  Strict interlacing
provides other rational symmetrizable realizations of the target pair by
Lemma~\ref{lem:realization}; the corollary prohibits reaching it while keeping
the particular form $H_K$ and the particular first compression fixed.  By
Dirichlet's theorem, the two residue classes used above contain infinitely
many primes.

\section{A complete height-one target slice}

Fix $K\geq5$, set $T=T_K$, and use the canonical update plane
\begin{equation}\label{eq:canonicalplane}
 U=(Te,T^2e),\qquad p_1=X,\quad p_2=X^2.
\end{equation}
Consider the target slice
\begin{equation}\label{eq:targetslice}
 f_{a,m}=X(X^2-aX-m),\qquad P_*=(X-2)(X+1),\qquad a,m\in\Q.
\end{equation}
For this plane direct cofactor calculation gives
\begin{align}
 \Phi_C&=\frac{(K-1)(2K-5)}{K^2}X,\label{eq:PhiCexplicit}\\
 \Phi_D&=\frac{2(K-1)(2K-5)}{K^3},\label{eq:PhiDexplicit}\\
 \Psi&=\frac{X^2(K-1)(2K-5)((K-2)X-1)}{K^3}.
 \label{eq:Psiexplicit}
\end{align}

\begin{theorem}[Exact rational image of the slice]\label{thm:slice}
For $K\geq5$, a rational symmetric rank-at-most-two update in the plane
\eqref{eq:canonicalplane} has output \eqref{eq:targetslice} if and only if
\begin{equation}\label{eq:slicecriterion}
 m=K,\qquad
 K(K-1)(2K-5)(K-1-a)(2a+K-4)\in\Q^2,
\end{equation}
where $\Q^2=\{q^2:q\in\Q\}$ includes zero.
For $K>4$, the target has strict Salem interlacing precisely when
\begin{equation}\label{eq:salemchamber}
 \frac{K-4}{2}<a<K-1.
\end{equation}
\end{theorem}

\begin{proof}
Write $S=\left(\begin{smallmatrix}r&s\\s&u\end{smallmatrix}\right)$ and
$d=ru-s^2$.  Substitute \eqref{eq:PhiCexplicit}--\eqref{eq:Psiexplicit}
and the cofactor matrices from \eqref{eq:matrixC}--\eqref{eq:matrixD} into
\eqref{eq:ranktwooutputs}.  Lemma~\ref{lem:sliceelimination} records the
coefficient solution explicitly.  In particular, coefficient comparison
forces
\begin{equation}\label{eq:mcompat}
 m=K.
\end{equation}
The remaining determinant condition is a genuine quadratic in $d$ whose
discriminant is
\begin{equation}\label{eq:discriminant}
 -16K^5(K-1)^3(2K-5)^3(a-K+1)(2a+K-4).
\end{equation}
Removing the evident square factor $16K^4(K-1)^2(2K-5)^2$ gives exactly the
square class in \eqref{eq:slicecriterion}.  The coefficient of $d^2$ is
$-(K-2)^2(K-1)^2(2K-5)^2\ne0$.  Conversely, if this class is a square
(including zero), the quadratic has a rational root $d$, back-substitution
gives a rational symmetric $S$, and Theorem~\ref{thm:ranktwoiff} supplies the
desired update.  This proves the classification.

For $m=K$, put $q(X)=X^2-aX-K$.  Since $q(0)<0$, its negative root lies in
$(-2,-1)$ precisely when $q(-2)>0$ and $q(-1)<0$, namely
\[
 2a+4-K>0,\qquad 1+a-K<0.
\]
These are exactly the two inequalities in \eqref{eq:salemchamber}.  They also
imply $q(2)<0$ for $K>4$, so the positive root lies above $2$.  Thus the
compressed roots $-1,2$ strictly interlace the three full roots.
\end{proof}

The eliminant computation in the proof is elementary but long.  The source
package includes a symbolic certificate that derives \eqref{eq:mcompat} and
\eqref{eq:discriminant} from the cofactor formulas without floating-point
arithmetic.

\begin{proposition}[Parametrization, finite square sieve, and forced signature]
\label{prop:sliceparam}
Fix an integer $K\geq5$, put
\[
 \Gamma_K=K(K-1)(2K-5),\qquad n=K-1-a,
\]
and restrict to $m=K$.  The rational image condition of
Theorem~\ref{thm:slice} is
\begin{equation}\label{eq:sliceconic}
 Y^2=\Gamma_K n\bigl(3(K-2)-2n\bigr).
\end{equation}
Apart from $(n,Y)=(0,0)$, all rational points of
\eqref{eq:sliceconic} are parametrized by $\ell\in\Q$ as
\begin{equation}\label{eq:sliceparam}
 n=\frac{3\Gamma_K(K-2)}{\ell^2+2\Gamma_K},\qquad
 a=K-1-n,\qquad Y=\ell n.
\end{equation}
If $a\in\Z$, strict Salem interlacing is equivalent to
\begin{equation}\label{eq:integersieve}
 1\leq n\leq\left\lfloor\frac{K+1}{2}\right\rfloor.
\end{equation}
Consequently, the strict integral image at fixed $K$ is obtained by the
finite test of retaining exactly those $n$ in \eqref{eq:integersieve} for
which the right side of \eqref{eq:sliceconic} is an integer square.

Every rational realization in the strict chamber with nonzero determinant
has $\det S<0$.  For a strict integral target every realizing branch has
nonzero determinant and is therefore a genuine indefinite rank-two update.
\end{proposition}

\begin{proof}
After substituting $a=K-1-n$, the square class in
\eqref{eq:slicecriterion} becomes the right side of
\eqref{eq:sliceconic}.  The line $Y=\ell n$ through $(0,0)$ gives
\eqref{eq:sliceparam}, and every point with $n\ne0$ has $\ell=Y/n$.
The inequalities \eqref{eq:salemchamber} become
$0<n<(K+2)/2$, which for integral $n$ is precisely
\eqref{eq:integersieve}.  An integer which is a square in $\Q$ is an
integer square.

In the normalization of Lemma~\ref{lem:sliceelimination}, the determinant
eliminant after $a=K-1-n$ may be chosen as
\begin{align}\label{eq:deteliminantn}
 \mathcal Q_{K,n}(d)={}&
 -(K-2)^2(K-1)^2(2K-5)^2d^2\\
 &-2K^2(K-1)(2K-5)\notag\\
 &\qquad{}\times
 \Bigl((2K^2-7K+2)n+3(K-1)(K-2)\Bigr)d\notag\\
 &-K^4\bigl((2K+1)n-3K+3\bigr)^2.\notag
\end{align}
For $K\geq5$ and $n>0$, the coefficient in parentheses on the second
line is positive.  Hence the sum of the two roots is negative and their
product is nonnegative.  When the square condition holds the roots are
rational, so every nonzero root is negative.  In the integral strict range
the constant term cannot vanish: $(2K+1)n=3(K-1)$ would force $K=4$ if
$n=1$, while for $n\geq2$ the left side exceeds the right side.  Both
determinant branches are therefore strictly negative.
\end{proof}

\begin{corollary}[Local--global square test on the exact slice]
\label{cor:slice-Hasse}
Put
\[
 D(K,a)=K(K-1)(2K-5)(K-1-a)(2a+K-4).
\]
The canonical fixed-plane surgery exists over $\Q$ if and only if
$D(K,a)\in\Q^2$.  Equivalently,
\[
 D(K,a)\geq0
 \quad\text{and}\quad
 D(K,a)\in\Q_p^2\quad\text{for every prime }p.
\]
Only the real place and the primes dividing the numerator or denominator of
$2D(K,a)$ need be tested.
\end{corollary}

\begin{proof}
The first statement is Theorem~\ref{thm:slice}.  A rational number is a
rational square exactly when it is nonnegative at the real place and a square
in every $\Q_p$.  For the finite assertion, write $D=A/B$ in lowest terms.
If the indicated local tests pass, then every prime valuation of $A$ and
$B$ is even and the real sign is positive; hence both $A$ and $B$ are
squares and $D\in\Q^2$.  The converse is immediate.  This is the elementary
one-dimensional case of the usual local square criterion; compare
\cite{cassels}.
\end{proof}

\begin{example}[One-prime all-depth solvability does not descend]
\label{ex:13-adic-no-descent}
For $K=6$, $a=2$, the target is
\[
 f_*=X(X^2-2X-6),\qquad P_*=(X-2)(X+1).
\]
The same-form norm equation is
\[
 (233y)^2+420z^2=41940,
\]
and has no point over $\Q_7$; hence no rational $H_6$-self-adjoint surgery
of any rank produces this pair.  In the canonical rank-two plane, however,
coefficient elimination gives
\[
\begin{aligned}
 r&=-\frac{3(1925d+11036)}{70},&
 s&=\frac{485d+2808}{30},&
 u&=-\frac{665d+3888}{210},
\end{aligned}
\]
and
\begin{equation}\label{eq:13adicF}
 F(d)=1225d^2+24570d+46656=0.
\end{equation}
Modulo $13$, a unit multiple of $F$ is $4d^2+3$.  The root $d=3$ is
simple because $8d\equiv11\pmod{13}$.  Hensel's lemma supplies a compatible
$d\in\Z_{13}$ and hence a canonical-plane solution over $\Z_{13}$, reducing
at the first level to
\[
 S\equiv\begin{pmatrix}4&3\\3&3\end{pmatrix}\pmod{13}.
\]
Thus solutions exist modulo every $13^r$, although no rational solution
exists.  Consistently,
\[
 \disc F=315^2\cdot3780,
 \qquad D(6,2)=3780\notin\Q^2.
\]
This example separates one-prime infinite-depth lifting from a genuine
global local--global test.
\end{example}

\section{The canonical near-shadow obstruction}

We next ask whether the exact slice can preserve the exterior trace up to a
vanishing error.  The answer is negative, for an elliptic reason.

\begin{theorem}[Arithmetic exclusion of the canonical near-shadow]
\label{thm:nearshadow}
For an integer $K\geq5$, the formally closest integral member of the target
slice is obtained by taking $a=K-2$.  It is a genuine height-one quartic
Salem target, but it is not in the rational image of the canonical plane.
More precisely, its square condition is
\begin{equation}\label{eq:nearquartic}
 Y^2=K(K-1)(2K-5)(3K-8),
\end{equation}
whose affine rational solutions are exactly
\[
 (0,0),(1,0),\left(\frac52,0\right),
 \left(\frac83,0\right),(2,\pm2),(4,\pm12).
\]
In particular, none has $K\geq5$.
\end{theorem}

\begin{proof}
Substitution of $a=K-2$ in \eqref{eq:slicecriterion} gives
\eqref{eq:nearquartic}.  Put
\begin{equation}\label{eq:near-map}
 z=24-\frac{40}{K},\qquad w=\frac{40Y}{K^2}.
\end{equation}
Away from $K=0$ this is birational, with inverse
\[
 K=\frac{40}{24-z},\qquad Y=\frac{40w}{(24-z)^2},
\]
and transforms \eqref{eq:nearquartic} into
\begin{equation}\label{eq:nearE}
 E:\quad w^2=(z+16)(z-8)(z-9).
\end{equation}
The point $(K,Y)=(0,0)$ extends to the point at infinity of $E$.

For completeness, we record a short $2$-isogeny descent.  Set $x=z-9$.
Then
\[
 E_1:\quad w^2=x(x+1)(x+25)=x^3+26x^2+25x,
\]
and the curve $2$-isogenous through $(0,0)$ is
\[
 E_1':\quad v^2=u^3-52u^2+576u=u(u-16)(u-36).
\]
For the standard descent maps $\alpha,\alpha'$, one has
\begin{equation}\label{eq:isogenyrank}
 2^{\rank E_1(\Q)}
 =\frac{|\alpha(E_1(\Q))|\,|\alpha'(E_1'(\Q))|}{4};
\end{equation}
see \cite[Chapter~III]{cremona}.  The first image is contained in the four
square classes $\{1,-1,5,-5\}$, all attained at
$O$, $(-1,0)$, $(5,30)$, and $(-5,20)$.

We claim that $\alpha'(E_1'(\Q))=\{1\}$.  Since
$u(u-16)(u-36)<0$ for $u<0$, only the positive square classes
$1,2,3,6$ can occur.  If $c\in\{2,3,6\}$ occurred, the homogeneous-space
criterion would give coprime integers $M,e$ and an integer $N$ satisfying
\begin{equation}\label{eq:descentquartic}
 N^2=cM^4-52M^2e^2+\frac{576}{c}e^4.
\end{equation}
For $c=2$, reduction modulo $4$ forces $M=2m$ and $e$ odd, and then
$N=4n$ with
\[
 n^2=2m^4-13m^2e^2+18e^4.
\]
Modulo $8$ the right side is $7$ if $m$ is odd and $2$ or $6$ if $m$ is
even, never a square.  For $c=3$ or $6$, reduction modulo $3$ forces
$3\mid Me$.  Coprimality leaves exactly one of $M,e$ divisible by $3$;
modulo $9$, \eqref{eq:descentquartic} then gives respectively
$N^2\equiv3$ or $6\pmod9$, again impossible.  Thus
$|\alpha'|=1$, and \eqref{eq:isogenyrank} gives $\rank E_1(\Q)=0$.

The eight visible rational points are
\[
 O,\ (-25,0),\ (-1,0),\ (0,0),\ (-5,\pm20),\ (5,\pm30).
\]
Both $7$ and $11$ are primes of good reduction, and direct counting gives
$\#E_1(\mathbf F_7)=\#E_1(\mathbf F_{11})=8$.  For every prime $\ell$, the
$\ell$-primary rational torsion injects at one of these two reductions;
hence $\#E_1(\Q)_{\rm tors}$ divides $\gcd(8,8)=8$.  The displayed eight
points are therefore all of $E_1(\Q)$.  This is the curve of Cremona label
$240a3$ (modern
LMFDB label $240.c2$); the same rank and torsion determination appears in
\cite[Proposition~5]{gonzalezxarles}.  Applying the inverse birational map
gives the asserted solutions of \eqref{eq:nearquartic}.

Finally, for $a=K-2$, the trace polynomial
$q_K=X^2-(K-2)X-K$ has one root in $(-2,-1)$ and one above $2$.
Its discriminant $K^2+4$ lies strictly between $K^2$ and $(K+1)^2$, while
$q_K(2)q_K(-2)=K(8-3K)<0$.  Thus its reciprocal lift is irreducible and
Salem, and the primitive interlacer $-(X+1)$ has height one.
\end{proof}

\begin{corollary}[No vanishing exterior shift in the integral slice]
\label{cor:novanishingshift}
Let
\[
 y_K=\frac{K-1+\sqrt{(K-1)^2+4}}2,
 \qquad y_{a,K}=\frac{a+\sqrt{a^2+4K}}2.
\]
Along any sequence of strictly interlacing integral targets in the rational
image of the canonical plane with $K\to\infty$,
\[
 \liminf |y_{a,K}-y_K|\geq\frac12.
\]
In particular, this complete height-one slice contains no $o(1)$ exterior
shadow surgery.
\end{corollary}

\begin{proof}
For the unique nearest integer $a_0=K-2$,
\[
 y_{a_0,K}=\frac{K-2+\sqrt{K^2+4}}2,
 \qquad
 y_{a_0,K}-y_K=-K^{-2}-K^{-3}+2K^{-4}+O(K^{-5}).
\]
Theorem~\ref{thm:nearshadow} excludes $a_0$ from the rational image for
every $K\geq5$.  In the strict chamber $a>0$ and
\[
 \frac{\partial y_{a,K}}{\partial a}
 =\frac12\left(1+\frac{a}{\sqrt{a^2+4K}}\right)>\frac12.
\]
Every other integral $a$ satisfies $|a-a_0|\geq1$, so the mean-value theorem
and the triangle inequality give
\[
 |y_{a,K}-y_K|>\frac12-|y_{a_0,K}-y_K|=\frac12-o(1).
\]
\end{proof}

\begin{corollary}[Fixed bounded shifts are finite]
\label{cor:boundedshift}
Fix an integer $h\leq0$ and put $a=K-2+h$.  There are only finitely many
integers $K\geq5$ for which this target lies in the rational image of the
canonical plane.  Consequently, no infinite integral image family with
bounded $h$ exists.
\end{corollary}

\begin{proof}
The square condition becomes
\[
 Y^2=(1-h)K(K-1)(2K-5)(3K+2h-8).
\]
For fixed $h\leq0$ the four roots on the right are distinct, so this is a
nonsingular affine genus-one curve.  Siegel's theorem gives only finitely
many integral points; see \cite[Chapter~IX]{silverman}.  A bounded set of
integer shifts is a finite union of such cases.
\end{proof}

\section{Exact Salem height collapse in rank two}

We first record a small sporadic certificate, then give an infinite family.

\begin{proposition}[The $K=5$ certificate]\label{prop:K5}
For $K=5$ and the plane \eqref{eq:canonicalplane}, take
\[
 S=\begin{pmatrix}-159/2&21\\21&-11/2\end{pmatrix},\qquad
 \det S=-\frac{15}{4}.
\]
Then
\begin{equation}\label{eq:K5outputs}
 f_S=X(X^2-X-5),\qquad P_S=(X-2)(X+1).
\end{equation}
The reciprocal lift
\begin{equation}\label{eq:K5salem}
 z^4-z^3-3z^2-z+1
\end{equation}
is an irreducible Salem polynomial, the update has rank exactly two, and the
output normal height is one.  This output is impossible in rank one.
\end{proposition}

\begin{proof}
Substitution in \eqref{eq:ranktwooutputs} gives \eqref{eq:K5outputs}.  The
trace roots are $(1\pm\sqrt{21})/2$, so
\[
 -2<\frac{1-\sqrt{21}}2<-1<0<2<\frac{1+\sqrt{21}}2.
\]
Irreducibility follows either from the quadratic trace polynomial and the
sign of the norm of $y^2-4$, or from the quartic criterion.  The columns of
$U$ are independent and $S$ is nonsingular.  The compressed interlacer
$-(X+1)$ is primitive, giving height one.  Finally
\[
 f_KP_S-f_SP_K=-\frac15X^2(X-2)(11X-1),
\]
which has only one factor $X-2$, contrary to the rank-one divisibility
\eqref{eq:pin-square-div}.
\end{proof}

\begin{theorem}[Infinite exact rank-two Salem surgery]\label{thm:infinite}
Let
\begin{equation}\label{eq:tconditions}
 \nu\geq3,\qquad \nu\equiv3\ \text{or }7\pmod {10},\qquad
 K=\nu^2+1,\qquad a=\frac{7\nu^2-3}{10}.
\end{equation}
There exists a rational symmetric $S_\nu$ in the fixed plane
\eqref{eq:canonicalplane} such that the update has rank exactly two and
\begin{equation}\label{eq:infiniteoutputs}
 f_\nu=X(X^2-aX-K),\qquad P_\nu=(X-2)(X+1).
\end{equation}
One branch satisfies
\begin{equation}\label{eq:detSt}
 \det S_\nu=-\frac{3(\nu-1)^2K^2}{10\nu^2(\nu+1)^2}<0.
\end{equation}
The reciprocal polynomial
\begin{equation}\label{eq:Rt}
 R_\nu(z)=z^4-az^3+(2-K)z^2-az+1
\end{equation}
is an irreducible Salem polynomial.  Moreover
\begin{equation}\label{eq:heightcollapse}
 \kappa(B_K)=K=\nu^2+1\longrightarrow\infty,\qquad
 \kappa(R_\nu)=1,
\end{equation}
and the same full--minor output cannot be produced by a rank-one update of
the base pencil.
\end{theorem}

\begin{proof}
The congruence in \eqref{eq:tconditions} makes $K,a$ integers.  Substitution
in the square class of Theorem~\ref{thm:slice} gives
\[
 K(K-1)(2K-5)(K-1-a)(2a+K-4)
 =\left(\frac35K(2K-5)\nu\right)^2.
\]
Hence a rational $S_\nu$ exists.  Solving the determinant quadratic with one
choice of square root gives \eqref{eq:detSt}; it is nonzero, so the update has
rank two.

Let $q=X^2-aX-K$.  Direct evaluation gives
\begin{align}\label{eq:qsigns}
 q(-2)&=2+\frac25K>0,&q(-1)&=-\frac3{10}K<0,\\
 q(0)&=-K<0,&q(2)&=6-\frac{12}{5}K<0.
\end{align}
Thus its roots satisfy
$-2<\alpha<-1<0<2<y$.  If $q$ were reducible over $\Q$, its root in
$(-2,-1)$ would be an integer, impossible.  Finally
\begin{equation}\label{eq:normsquare}
 \N_{\Q(y)/\Q}(y^2-4)=q(2)q(-2)
 =-\frac{12}{25}(2K-5)(K+5)<0.
\end{equation}
A square in the real quadratic trace field has nonnegative norm.  Hence
$y^2-4$ is not a square, and the reciprocal lift is irreducible.  The root
placement proves the Salem signature.

The output interlacer $-(X+1)$ is primitive, so $\kappa(R_\nu)=1$, whereas
Proposition~\ref{prop:exactheight} gives the base height.  Finally
\begin{equation}\label{eq:crossinfinite}
 f_KP_\nu-f_\nu P_K=
 \frac1KX^2(X-2)
 \bigl[(aK-K^2+2K-1)X+a\bigr].
\end{equation}
At $X=2$ the bracket equals
$-\frac3{10}(2K-5)(K-2)\ne0$.  Thus $X-2$ occurs only once, violating the
rank-one square divisibility.
\end{proof}

The first two instances are
\[
\begin{array}{c@{\quad}c@{\quad}c@{\quad}l}
\toprule
\nu&K&a&R_\nu(z)\\
\midrule
3&10&6&z^4-6z^3-8z^2-6z+1,\\
7&50&34&z^4-34z^3-48z^2-34z+1.\\
\bottomrule
\end{array}
\]

\section{Scope, novelty, and open directions}

Theorems~\ref{thm:rankone-iff} and \ref{thm:ranktwoiff} classify the rational
image of a \emph{fixed cyclic realization} together with its fixed first
principal compression.  They are not claims about arbitrary matrices,
arbitrary symmetrizers, or the full Jacobi inverse problem.  Rank-$m$
characteristic-polynomial assignment, restricted-rank eigenvalue
modification, and algebraic relations among principal minors have substantial
independent literatures; see
\cite{merzelm,arbenzgangergolub,alahmadiehvinzant}.  Within that landscape,
the feature isolated here is the paired full--compression assignment law for
one cyclic self-adjoint model, including a converse, followed by exact
arithmetic and Salem applications.  We do not claim novelty for the secular,
Rayleigh-square, local--global, elliptic-rank, or quartic-Salem ingredients
individually; the new statements are their constrained simultaneous-surgery
consequences.

The negative results close four sharply delimited routes.  First, pointwise
full and compressed eigenspace preservation imposes the rank bound of
Theorem~\ref{thm:allrankbarrier}.  Second, semidefinite perturbations meeting
the two ordered output pins and strict integral Salem signature are excluded
in every rank by Theorem~\ref{thm:allrank-semidefinite}.  Third, the binary
norm condition in Theorem~\ref{thm:sameH-norm} and its infinite local
obstruction apply at every rank while the particular canonical form $H_K$
is retained.  Fourth, in the particular plane $(T_Ke,T_K^2e)$, the exact
slice admits neither a vanishing integral exterior shift nor an infinite
family with any fixed bounded shift.  The rational-point determination used
for the nearest shift is known for the elliptic curve $240a3$
\cite{gonzalezxarles}; the new content here is its derivation as the exact
spectral feasibility curve and the resulting surgery obstruction.

The positive results are complementary.  The $K=5$ certificate and
Theorem~\ref{thm:infinite} show that the rank-one square obstruction is
genuinely removed and that exact arithmetic height collapse occurs in rank
two.  Proposition~\ref{prop:sliceparam} also shows that every strict branch
in the complete canonical slice is forced to be indefinite, rather than this
being an accident of the exhibited family.

None of these statements excludes a construction that changes $H_K$, uses a
different rational two-plane, increases the ambient order, or gives up the
paired eigenspace pins.  Nor do the successful updates preserve tridiagonal
Jacobi sparsity or give an output for every integer $K$.  The infinite family
has an exterior shift of order $K$, so it does not prove a Boyd closure
conjecture.  A closure-relevant construction would still have to combine
higher order, simultaneous integrality, bounded or vanishing exterior shift,
Salem signature, and irreducibility.  These are genuine open directions, not
consequences claimed in this paper.

\appendix

\section{The target-slice elimination}\label{app:sliceelimination}

The following calculation makes the elimination in
Theorem~\ref{thm:slice} reproducible without relying on a computer-algebra
black box.

\begin{lemma}[Affine coefficient solution]\label{lem:sliceelimination}
Put $\delta=(K-1)(2K-5)$ and write
$S=\left(\begin{smallmatrix}r&s\\s&u\end{smallmatrix}\right)$ and
$d=\det S$.  Coefficient comparison in \eqref{eq:ranktwooutputs} with target
\eqref{eq:targetslice} first forces $m=K$ and then gives
\begin{align*}
r={}&\frac{aK^4-4aK^3+4aK^2-3aK-K^6+5K^5-9K^4+9K^3-5K^2+K}
 {\delta}\\
&+\frac{-K^4+5K^3-9K^2+8K-3}{K}\,d,\\
s={}&\frac{-2aK^3+4aK^2-aK+2K^5-8K^4+10K^3-6K^2+2K}
 {2\delta}\\
&+\frac{2K^3-8K^2+9K-4}{2K}\,d,\\
u={}&\frac{aK^2-K^4+3K^3-2K^2}{\delta}
 +\frac{-K^2+3K-1}{K}\,d.
\end{align*}
After clearing denominators, let $\mathcal Q_{K,a}(d)$ be the numerator of
$ru-s^2-d$.  It is a quadratic with leading coefficient
\[
 -(K-2)^2(K-1)^2(2K-5)^2
\]
and discriminant
\[
 -16K^5(K-1)^3(2K-5)^3
 (a-K+1)(2a+K-4).
\]
Consequently $\mathcal Q_{K,a}$ has a rational root if and only if the
square-class condition \eqref{eq:slicecriterion} holds.
\end{lemma}

\begin{proof}
Insert $p_1=X,p_2=X^2$ in the remainder formulas of
Lemma~\ref{lem:rankone-normal}, polarize to obtain $C$ and $D$, and compare
the coefficients of $f_S-f_{a,m}$ and $P_S-P_*$.  Gaussian elimination over
$\Q(K,a,d)$ gives the three displayed affine expressions and $m=K$.
Substitution in $ru-s^2=d$ gives $\mathcal Q_{K,a}$; expanding its leading
coefficient and its discriminant gives the displayed factors.  Its leading
coefficient is nonzero for $K\ge5$.  Finally its discriminant equals
\[
 \bigl(4K^2(K-1)(2K-5)\bigr)^2
 K(K-1)(2K-5)(K-1-a)(2a+K-4),
\]
which proves the last assertion, including the zero-discriminant case.
\end{proof}

\section{Explicit matrix certificate for the infinite family}
\label{app:explicit}

For completeness we display the branch used in
Theorem~\ref{thm:infinite}.  Put $K=\nu^2+1$ and
$D=10\nu^2(\nu+1)^2$.  Define
\begin{align*}
 U_{11}={}&2\nu^{10}+16\nu^9+9\nu^8+2\nu^7+13\nu^6+12\nu^5
 +10\nu^4-4\nu^3+6\nu^2+4\nu+2,\\
 U_{12}={}&2\nu^8+16\nu^7+9\nu^6+2\nu^5+11\nu^4+2\nu^3
 +6\nu^2-2\nu+2,\\
 U_{22}={}&2\nu^6+16\nu^5+9\nu^4+2\nu^3+9\nu^2-8\nu+2.
\end{align*}
Then one may take
\begin{equation}\label{eq:explicitSnu}
 S_\nu=\frac K D
 \begin{pmatrix}-U_{11}&U_{12}\\U_{12}&-U_{22}\end{pmatrix}.
\end{equation}
Direct expansion gives \eqref{eq:detSt} and the output pair
\eqref{eq:infiniteoutputs}.

An entirely rational base realization is also explicit.  With
$\mathcal B_K,b_1,b_2$ as in \eqref{eq:canonical-b}, set
\begin{align*}
 a_0&=\frac{K^2-3K+1}{K},&
 a_1&=\frac{4K^4-12K^3+6K^2-3K+1}{K\mathcal B_K},\\
 a_2&=-\frac{2(K^2-3K+1)}{\mathcal B_K}.&&
\end{align*}
With
\[
 T_K=\begin{pmatrix}a_0&1&0\\b_1&a_1&1\\0&b_2&a_2\end{pmatrix},
 \qquad H_K=\operatorname{diag}\left(1,b_1^{-1},(b_1b_2)^{-1}\right),
\]
one has $H_KT_K=T_K^{\mathsf T}H_K$, the characteristic pair is
$(f_K,P_K)$, and the update plane is exactly
$U_K=(T_Ke,T_K^2e)$.  Substitution of \eqref{eq:explicitSnu} verifies the
two target characteristic polynomials over $\Q(\nu)$.

\section{Exact certificates}

All displayed elimination identities and the two explicit Salem surgeries
were verified over the rational function field with exact symbolic
arithmetic.  The companion scripts included with the source package check:
\begin{enumerate}[label=\textup{(\arabic*)}]
\item the rank-one coefficient elimination and all $K=5,\ldots,150$ against
an enlarged integral search window;
\item the cofactor and Gram--Pl\"ucker identities;
\item the $K=5$ rational $H$-self-adjoint matrix realization;
\item the target-slice eliminant and discriminant;
\item the same-form norm identity and the first local obstruction instance;
\item the parametrized slice conic, determinant eliminant, and forced
negative determinant;
\item the birational near-shadow identity, the displayed rational points,
and the point counts over $\mathbf F_7$ and $\mathbf F_{11}$;
\item and the infinite family, including \eqref{eq:detSt},
\eqref{eq:crossinfinite}, rank, and reciprocal output.
\end{enumerate}
These computations certify displayed algebra; they do not replace the proofs
of the converse or the root-location arguments.

\begin{thebibliography}{99}

\bibitem{alahmadiehvinzant}
A. Al Ahmadieh and C. Vinzant,
\emph{Characterizing principal minors of symmetric matrices via
determinantal multiaffine polynomials},
J. Algebra \textbf{638} (2024), 255--278,
\href{https://doi.org/10.1016/j.jalgebra.2023.09.030}
{doi:10.1016/j.jalgebra.2023.09.030}.

\bibitem{arbenzgangergolub}
P. Arbenz, W. Gander, and G. H. Golub,
\emph{Restricted rank modification of the symmetric eigenvalue problem:
theoretical considerations},
Linear Algebra Appl. \textbf{104} (1988), 75--95,
\href{https://doi.org/10.1016/0024-3795(88)90309-6}
{doi:10.1016/0024-3795(88)90309-6}.

\bibitem{bhatia}
R. Bhatia,
\emph{Matrix Analysis},
Graduate Texts in Mathematics 169, Springer, New York, 1997,
\href{https://doi.org/10.1007/978-1-4612-0653-8}
{doi:10.1007/978-1-4612-0653-8}.

\bibitem{boydquartic}
D. W. Boyd,
\emph{Salem numbers of degree four have periodic expansions},
in \emph{Th\'eorie des nombres (Quebec, 1987)}, de Gruyter, Berlin, 1989,
57--64.

\bibitem{cassels}
J. W. S. Cassels,
\emph{Rational Quadratic Forms},
London Mathematical Society Monographs 13, Academic Press, London, 1978.

\bibitem{cremona}
J. E. Cremona,
\emph{Algorithms for Modular Elliptic Curves},
2nd ed., Cambridge University Press, Cambridge, 1997.

\bibitem{delrio}
R. del Rio, M. Kudryavtsev, and L. O. Silva,
\emph{Inverse problems for Jacobi operators II: mass perturbations of
semi-infinite mass-spring systems},
\href{https://arxiv.org/abs/1106.4598}{arXiv:1106.4598}.

\bibitem{dptz}
P. B. Denton, S. J. Parke, T. Tao, and X. Zhang,
\emph{Eigenvectors from eigenvalues: a survey of a basic identity in linear
algebra},
Bull. Amer. Math. Soc. (N.S.) \textbf{59} (2022), 31--58,
\href{https://doi.org/10.1090/bull/1722}{doi:10.1090/bull/1722}.

\bibitem{deboor}
C. de Boor and G. H. Golub,
\emph{The numerically stable reconstruction of a Jacobi matrix from spectral
data},
Linear Algebra Appl. \textbf{21} (1978), 245--260,
\href{https://doi.org/10.1016/0024-3795(78)90086-1}
{doi:10.1016/0024-3795(78)90086-1}.

\bibitem{hochstadt}
H. Hochstadt,
\emph{On the construction of a Jacobi matrix from spectral data},
Linear Algebra Appl. \textbf{8} (1974), 435--446,
\href{https://doi.org/10.1016/0024-3795(74)90077-9}
{doi:10.1016/0024-3795(74)90077-9}.

\bibitem{kula}
A. Kula, M. Wojtylak, and J. Wysocza\'nski,
\emph{Rank two perturbations of matrices and operators and operator model for
$t$-transformation of probability measures},
J. Funct. Anal. \textbf{272} (2017), 1147--1181,
\href{https://doi.org/10.1016/j.jfa.2016.10.021}
{doi:10.1016/j.jfa.2016.10.021}.

\bibitem{gonzalezxarles}
E. Gonz\'alez-Jim\'enez and X. Xarles,
\emph{On a conjecture of Rudin on squares in arithmetic progressions},
LMS J. Comput. Math. \textbf{17} (2014), no.~1, 58--76,
\href{https://doi.org/10.1112/S1461157013000259}
{doi:10.1112/S1461157013000259}.

\bibitem{merzel}
J. Merzel, J. Min\'a\v c, L. Muller, F. W. Pasini, and T. T. Nguyen,
\emph{Spectral perturbation by rank one matrices},
\href{https://arxiv.org/abs/2111.10624}{arXiv:2111.10624}.

\bibitem{merzelm}
J. L. Merzel, J. Min\'a\v c, T. T. Nguyen, and F. W. Pasini,
\emph{Spectral perturbation by rank-$m$ matrices},
Operators and Matrices \textbf{17} (2023), no.~3, 867--874,
\href{https://doi.org/10.7153/oam-2023-17-58}
{doi:10.7153/oam-2023-17-58}.

\bibitem{silverman}
J. H. Silverman,
\emph{The Arithmetic of Elliptic Curves},
2nd ed., Graduate Texts in Mathematics 106, Springer, New York, 2009,
\href{https://doi.org/10.1007/978-0-387-09494-6}
{doi:10.1007/978-0-387-09494-6}.

\bibitem{zaimi}
T. Za\"imi,
\emph{Quartic Salem numbers which are Mahler measures of non-reciprocal
2-Pisot numbers},
J. Th\'eor. Nombres Bordeaux \textbf{32} (2020), no.~3, 877--889,
\href{https://doi.org/10.5802/jtnb.1145}{doi:10.5802/jtnb.1145}.

\end{thebibliography}

\end{document}
